\BeleskeID{05-s022}
% element: 05-s022:d01
Posmatrajmo najpre nenegativne apsolutne vrednosti $a,b$. Oduzimanje negativnog broja daje zbir; oduzimanje pozitivnog od negativnog daje negativan zbir njihovih apsolutnih vrednosti. Preostala dva slučaja mogu se zapisati kao sabiranje sa promenjenim znakom drugog operanda.

Kada je $a<b$, apsolutna vrednost negativnog rezultata je $b-a$. U $n$-bitnom drugom komplementu njegov zapis je $2^n-(b-a)$. Pregrupisavanjem dobijamo $a+(2^n-b)$: pozitivan prvi operand sabira se sa zapisom negativnog drugog operanda. Ekvivalentnost zapisa odnosi se na aritmetiku po modulu $2^n$; višak izvan registra ne čuva se u rezultatu.
